% !TEX program = lualatex
\documentclass[12pt, a4paper]{article}
\usepackage{master}
% tikz: 제1실체 의존 구조 방사 다이어그램용
\usepackage{tikz}
\usetikzlibrary{calc, positioning, arrows.meta}

% 저술 정의
\DeclareWorkGR{Cat}{범주론}{Κατηγορίαι}{Categoriae}
\DeclareWorkGR{DI}{명제론}{Περὶ ἑρμηνείας}{De interpretatione}
\DeclareWorkGR{Met}{형이상학}{Μεταφυσικά}{Metaphysica}
\DeclareWorkGR{GC}{생성소멸론}{Περὶ γενέσεως καὶ φθορᾶς}{De generatione et corruptione}
\DeclareWorkGR{Phys}{자연학}{Φυσικά}{Physica}

\fancyhead[L]{\footnotesize\color{midgray}오르가논 강독}
\fancyhead[R]{\footnotesize\color{midgray}제2주}

\begin{document}
\thispagestyle{firstpage}

\weeksection{제2주}{범주론 1a1--4a21}{선행 개념과 실체}

\subsection{범주론의 첫 마디: 동음이의, 동의, 파생어 (1장, 1a1--15)}

\subsubsection{왜 이 구별로 시작하는가}

범주론은 존재론적 분류 작업이지만, 그 첫 걸음은 의외로 의미론적(semantic)입니다. 아리스토텔레스는 사물을 분류하기 전에, 먼저 이름(\gr{ὄνομα})과 사물(\gr{πρᾶγμα})의 관계를 정리합니다. 같은 이름이 서로 다른 사물에 적용될 때, 그 적용 방식이 같은가 다른가를 구별하는 것이 출발점입니다. 이것은 단순한 언어학적 예비 작업이 아니라, 범주적 분류 자체의 전제 조건입니다 --- 이름이 같다고 해서 사물도 같다고 가정하면, 범주의 구별 자체가 불가능해지기 때문입니다.

범주론 1장과 이후 장들의 연결 관계는 고대부터 논의되어 왔습니다. 근래에 마이클 프레데(Michael Frede)는 1장의 동음이의·동의 구별이 2장의 `주어에 대해 말해지는 것'(said-of)과 `주어 안에 있는 것'(present-in) 관계와 어떻게 연결되는지를 분석하여, 1장이 단순한 예비 작업이 아니라 범주론 전체의 존재론적 틀을 준비하는 것이라고 주장했습니다.\footnote{Michael Frede, ``Categories in Aristotle,'' in \gri{Essays in Ancient Philosophy} (Minneapolis: University of Minnesota Press, 1987), 29--48. 프레데의 이 논문은 범주론 해석사에서 전환점이 된 것으로 평가됩니다.} 이 점은 1장을 읽을 때 염두에 둘 가치가 있습니다.

\subsubsection{동음이의적인 것들 (\gr{ὁμώνυμα})}

\srcquote{\gr{Ὁμώνυμα λέγεται ὧν ὄνομα μόνον κοινόν, ὁ δὲ κατὰ τοὔνομα λόγος τῆς οὐσίας ἕτερος, οἷον ζῷον ὅ τε ἄνθρωπος καὶ τὸ γεγραμμένον.}}{`이름만 공유하고, 이름에 대응하는 존재의 정의(\gr{λόγος τῆς οὐσίας})가 다른 것들을 동음이의적이라 부른다. 예컨대 사람도 동물(\gr{ζῷον})이고 그림도 동물이다.'}{범주론 1a1--2}

아리스토텔레스의 예시에서 핵심은 그리스어 \gr{ζῷον}이 `살아 있는 것'이라는 뜻과 `(살아 있는 것의) 그림'이라는 뜻을 모두 가진다는 점입니다.\footnote{범주론 그리스어 텍스트: L.\ Minio-Paluello, ed., \gri{Aristotelis Categoriae et Liber de Interpretatione}, OCT (Oxford: Clarendon Press, 1949; repr.\ with corrections, 1956). Eng.\ trans.: J.\ L.\ Ackrill, \gri{Aristotle's Categories and De Interpretatione}, Clarendon Aristotle Series (Oxford: Clarendon Press, 1963). 애크릴은 1a1--6에 대한 주석에서 \gr{ζῷον}의 이중 의미를 상세히 논의합니다.} 사람에게 적용될 때와 그림에 적용될 때, `동물임'(\gr{τὸ ζῴῳ εἶναι})의 정의는 전혀 다릅니다.

여기서 반드시 짚어야 할 것은, 아리스토텔레스가 말하는 `동음이의'가 현대 언어학의 `동음이의어'(homonymy)와 정확히 같지 않다는 점입니다. 현대의 동음이의어는 단어의 속성이지만, 아리스토텔레스의 동음이의는 사물의 속성입니다 --- `동음이의적으로 불리는 것들'(\gr{ὁμώνυμα λέγεται})이지, `동음이의적인 단어'가 아닙니다. 야코 힌티카(Jaakko Hintikka)와 여러 학자들이 지적하듯이, 아리스토텔레스 자신도 다른 저작에서는 이 구별을 때때로 느슨하게 사용하여 단어 자체를 동음이의적이라고 부르기도 합니다(\citework{GC}{I.6, 322b29} 이하). 그러나 범주론에서의 엄밀한 정의는 분명히 사물에 관한 것입니다.\footnote{Jaakko Hintikka, ``Semantical Games, the Alleged Ambiguity of `Is' and Aristotelian Categories,'' \gri{Synthese} 54 (1983): 443--468.}

\subsubsection{동의적인 것들 (\gr{συνώνυμα})}

\srcquote{\gr{Συνώνυμα δὲ λέγεται ὧν τό τε ὄνομα κοινὸν καὶ ὁ κατὰ τοὔνομα λόγος τῆς οὐσίας ὁ αὐτός, οἷον ζῷον ὅ τε ἄνθρωπος καὶ ὁ βοῦς.}}{`이름도 공유하고, 이름에 대응하는 존재의 정의도 같은 것들을 동의적이라 부른다. 예컨대 사람도 동물이고 소도 동물이다.'}{범주론 1a6--8}

여기서는 같은 이름 `동물'이 사람과 소 모두에게 같은 의미로 적용됩니다 --- `동물임'의 정의가 양자에게 동일합니다. 현대적 용어법과의 차이에 다시 주목해야 합니다. 현대 언어학에서 `동의어'(synonym)는 다른 단어가 같은 의미를 갖는 것이지만, 아리스토텔레스의 동의(\gr{συνώνυμα})는 같은 이름이 같은 정의로 여러 사물에 적용되는 것입니다.

이 구별이 왜 철학적으로 중요한지를 보여주는 예를 하나 생각해 봅시다. `건강한'(\gr{ὑγιεινός})이라는 단어는 사람에게도, 음식에게도, 안색에게도 적용됩니다. 그런데 사람이 건강하다는 것과, 음식이 건강하다는 것과, 안색이 건강하다는 것의 정의는 같을까요, 다를까요? 사람은 건강을 소유하고, 음식은 건강을 생산하고, 안색은 건강을 표시합니다. 정의가 다르니 이것은 동음이의일까요? 그렇습니다. 그러나 완전히 아무 관련 없는 우연적 동음이의도 아닙니다 --- 세 경우 모두 건강이라는 하나의 핵심 개념(\gr{ἕν})에 대한 관계를 통해 `건강하다'고 불립니다.

아리스토텔레스는 \citework{Met}{\gr{Γ}권}에서 이 현상을 `하나에 관련된'(\gr{πρὸς ἕν}) 동음이의, 곧 `핵심에 관련된 동음이의'(core-dependent homonymy)라는 개념으로 발전시킵니다.\footnote{`핵심에 관련된 동음이의'라는 현대적 표현은 G.E.L.\ Owen, ``Logic and Metaphysics in Some Earlier Works of Aristotle,'' in \gri{Aristotle and Plato in the Mid-Fourth Century}, ed.\ I.\ Düring and G.E.L.\ Owen (Göteborg, 1960), 163--190에서 비롯됩니다. `focal meaning'이라는 용어가 널리 쓰이게 된 것은 이 논문의 영향입니다.} 스탠퍼드 철학 백과사전의 「아리스토텔레스의 범주론」 항목에서 스터트만이 설명하듯이, 아리스토텔레스에 따르면 존재하는 것들이 `존재자'(\gr{ὄντα})라고 불리는 것도 이와 같은 구조입니다 --- 실체, 양, 질 등은 어떤 하나의 류(존재자)에 속하기 때문에 `존재자'라고 불리는 것이 아니라, 모두 제1의 존재자, 곧 실체에 대한 어떤 관계를 맺기 때문에 `존재자'라고 불립니다(형이상학 1003a35 이하).\footnote{Paul Studtmann, ``Aristotle's Categories,'' in \gri{Stanford Encyclopedia of Philosophy}, §5.} 범주론의 이 서두는 그러한 후기 논의의 씨앗을 담고 있습니다.

\subsubsection{파생어 (\gr{παρώνυμα})}

아리스토텔레스는 세 번째로 파생어(\gr{παρώνυμα})를 도입합니다. 문법에서 이름을 얻되 어미가 달라지는 것들입니다 --- 문법(\gr{γραμματική})에서 문법가(\gr{γραμματικός})가, 용기(\gr{ἀνδρεία})에서 용감한 자(\gr{ἀνδρεῖος})가 나옵니다(1a12--15). 이 구별은 나중에 실체와 속성의 관계를 설명할 때 중요한 역할을 합니다. 속성이 실체 `안에 있을' 때, 그 실체는 속성에서 파생된 이름으로 불리는 경우가 많습니다 --- 몸(\gr{σῶμα}) 안에 흼(\gr{λευκότης})이 있으면, 몸은 `흰'(\gr{λευκόν})이라고 불립니다. 파생어 관계가 `주어 안에 있는 것' 관계의 언어적 표지(marker) 역할을 하는 셈입니다. 이것이 바로 프레데가 주장한 1장과 2장의 구조적 연결입니다 --- 동의(synonymy)는 `주어에 대해 말해지는 것'(said-of)에, 파생어(paronymy)와 동음이의(homonymy)는 `주어 안에 있는 것'(present-in)에 대응합니다.\footnote{Frede, ``Categories in Aristotle,'' 36--41.}

\subsection{결합과 비결합, 사중 분류 (2장, 1a16--1b9)}

\subsubsection{결합 표현과 비결합 표현}

2장의 서두(1a16--19)에서 아리스토텔레스는, 말해지는 것들(\gr{τὰ λεγόμενα}) 가운데 결합을 포함하는 것(`사람이 달린다', `사람이 이긴다')과 결합 없이 말해지는 것(`사람', `소', `달린다', `이긴다')을 구별합니다. 결합 없이 말해지는 것은 그 자체로 참도 거짓도 아닙니다 --- `사람'이라는 말 자체는 참도 거짓도 아니며, `달린다'도 마찬가지입니다. 참·거짓은 결합에서만 성립합니다. 이 구별은 명제론에서 본격적으로 전개될 주장, 곧 참·거짓을 가리는 것은 명제(\gr{ἀπόφανσις})뿐이라는 주장의 예고입니다.

\subsubsection{`주어 안에 있는 것'의 정의}

사중 분류를 이해하기 위해서는 먼저 아리스토텔레스가 `주어 안에 있는 것'(\gr{ἐν ὑποκειμένῳ})을 어떻게 정의하는지를 정확히 파악해야 합니다. 이 정의는 범주론 전체의 존재론적 구조를 받치는 토대입니다.

\srcquote{\gr{ἐν ὑποκειμένῳ δὲ λέγω ὃ ἔν τινι μὴ ὡς μέρος ὑπάρχον ἀδύνατον χωρὶς εἶναι τοῦ ἐν ᾧ ἐστίν.}}{`주어 안에 있다고 내가 말하는 것은, 어떤 것 안에 부분으로서가 아니라 존재하면서, 그것이 들어 있는 것과 분리되어 존재할 수 없는 것이다.'}{범주론 1a24--25}

여기서 아리스토텔레스는 두 가지 조건을 제시합니다. 첫째, 부분(\gr{μέρος})이 아닐 것 --- 소크라테스의 손은 소크라테스 `안에 있지만', 이 의미에서의 `안에 있는 것'이 아닙니다. 둘째, 분리 불가능할 것 --- 소크라테스 안에 있는 이 특정한 흼은 소크라테스 없이 존재할 수 없습니다.\footnote{`분리되어 존재할 수 없다'(\gr{ἀδύνατον χωρὶς εἶναι})는 조건의 정확한 의미는 논쟁의 대상입니다. 이것이 특정한 주어와의 분리를 뜻하는지, 아니면 어떤 주어든 주어 일반과의 분리를 뜻하는지가 문제입니다. 애크릴은 후자의 해석을 지지합니다(Ackrill, \gri{Categories}, 1a24--25 주석).}

마크 코언(Marc Cohen)은 이 두 관계를 다음과 같이 정리합니다. `주어에 대해 말해지는 것'(said-of)은 존재자를 보편자와 개별자로 나누는 관계이고, `주어 안에 있는 것'(present-in)은 존재자를 실체와 비실체로 나누는 관계입니다.\footnote{Marc Cohen, ``Predication and Ontology: Categories,'' 강의 자료, University of Washington.} 이 두 분할이 교차하면 네 가지 경우가 생깁니다.

\subsubsection{사중 분류}

2장의 핵심(1a20--1b9)은 존재하는 것들(\gr{τὰ ὄντα})을 두 가지 관계의 결합으로 분류하는 사중 분류입니다.

\needspace{8\baselineskip}
{%
\small
\setlength{\extrarowheight}{4pt}%
\renewcommand{\arraystretch}{1.3}%
\setlength{\tabcolsep}{5pt}%
\rowcolors{2}{altrowbg}{white}%
\noindent\centering
\begin{tabular}{|>{\centering\arraybackslash}m{30mm}|>{\centering\arraybackslash}m{45mm}|>{\centering\arraybackslash}m{45mm}|}
\hline
\rowcolor{tableheadblue}
{\color{h1navy}\bfseries} &
{\color{h1navy}\bfseries 주어 안에 있지 않음} &
{\color{h1navy}\bfseries 주어 안에 있음} \\ \hline
{\bfseries 주어에 대해\newline 말해짐} &
(a) 제2실체 \newline {\scriptsize 예: 사람, 동물} &
(c) 보편적 속성 \newline {\scriptsize 예: 지식, 색} \\ \hline
{\bfseries 주어에 대해\newline 말해지지 않음} &
\textbf{(d) 제1실체} \newline {\scriptsize 예: 개별 사람, 개별 말} &
(b) 개별적 속성 \newline {\scriptsize 예: 이 흼, 이 문법 지식} \\ \hline
\end{tabular}
\par
}

이 도표에서 아리스토텔레스가 제시하는 순서는 (a), (b), (c), (d)이지만, 논리적 핵심은 (d), 곧 제1실체입니다. 제1실체는 어떤 것에 대해서도 말해지지 않고 어떤 것 안에도 있지 않은 것 --- 두 관계 모두에서 `의존하는 쪽'이 아니라 `의존받는 쪽'인 궁극적 주어(\gr{ὑποκείμενον})입니다.

스탠퍼드 철학 백과사전의 「아리스토텔레스의 범주론」 항목에서 스터트만은, 이 사중 분류가 아리스토텔레스의 존재론을 뚜렷이 반플라톤적인 것으로 만든다고 지적합니다. 플라톤이 추상적인 형상(\gr{εἶδος}/\gr{ἰδέα})을 가장 실재적인 것으로 보았던 반면, 범주론은 구체적 개별자를 존재론적 토대로 놓기 때문입니다.\footnote{Paul Studtmann, ``Aristotle's Categories,'' in \gri{Stanford Encyclopedia of Philosophy}, §1.} 코언 역시 이 점에 주목하면서, 플라톤에게 술어화(predication)는 일반적으로 형상에의 참여(\gr{μέθεξις})로 설명되었지만, 아리스토텔레스는 이것이 지나치게 단순하다고 보았으며 본질적 술어(said-of)와 우연적 속성(present-in)을 엄격히 구별해야 한다고 주장했다고 설명합니다.

\subsubsection{개별적 속성의 문제}

사중 분류에서 가장 논쟁적인 칸은 (b), 곧 `주어 안에 있지만 주어에 대해 말해지지 않는 것' --- 개별적 속성(individual accident)입니다. 아리스토텔레스는 `개별적 문법 지식'(\gr{ἡ τὶς γραμματική})과 `개별적 흼'(\gr{τὸ τὶ λευκόν})을 예로 듭니다. 그런데 개별적인 흼이란 무엇일까요? 소크라테스의 흼과 플라톤의 흼은 수적으로 다른 두 개의 흼일까요?

이 물음은 마이클 웨딘(Michael V.\ Wedin)의 1993년 논문 이래로 활발히 논의되어 왔습니다.\footnote{Michael V.\ Wedin, ``Nonsubstantial Individuals,'' \gri{Phronesis} 38 (1993): 137--165.} 웨딘은 아리스토텔레스의 비실체적 개별자(non-substantial individual)가 기존의 두 가지 해석 --- 개별화된 속성(trope)으로 보는 해석과 단순히 보편자의 사례로 보는 해석 --- 어느 쪽으로도 완전히 포착되지 않는다고 주장했습니다. 이 논쟁의 핵심은 아리스토텔레스의 존재론이 보편자-개별자 구분을 실체 범주에만 한정하는가, 아니면 모든 범주에 걸쳐 적용하는가의 문제입니다. 사중 분류의 구조 자체는 후자를 함축합니다 --- 비실체 범주에도 보편적인 것(c)과 개별적인 것(b)이 있다는 것이니까요.

여기서 한 가지 사고 연습을 해봅시다. `주어에 대해 말해지는 것'과 `주어 안에 있는 것'은 얼핏 비슷해 보입니다 --- 둘 다 어떤 것이 어떤 것에 `속하는' 관계를 나타내는 것 아닌가? 그러나 아리스토텔레스는 이 둘을 엄격히 구별합니다. `사람이 소크라테스에 대해 말해진다'와 `흼이 소크라테스 안에 있다'는 어떻게 다를까요? 전자에서 소크라테스는 사람이라는 종(\gr{εἶδος})의 사례인 반면, 후자에서 흼은 소크라테스의 본질을 구성하지 않는 우연적 속성입니다. 이 구별은 본질적 술어와 우연적 속성의 차이, 나아가 `무엇임'(\gr{τί ἐστι})과 `어떠함'(\gr{ποῖόν τι})의 차이를 반영합니다. said-of 관계에서는 이름과 정의가 함께 전이되지만(동의적 술어화), present-in 관계에서는 이름은 전이될 수 있어도 정의는 전이되지 않습니다 --- 소크라테스를 `흰'이라고 부를 수 있지만, `흼의 정의'는 소크라테스에게 적용되지 않습니다(2a19--34).

\subsection{류 술어의 전이성 (3장, 1b10--24)}

3장은 짧지만 중요한 원리를 확립합니다. 어떤 것이 다른 것에 대해 주어로서 술어될 때, 그 술어에 대해 말해지는 모든 것이 원래의 주어에 대해서도 말해진다는 원리입니다(1b10--15). 예컨대 `사람'은 개별 사람에 대해 말해지고, `동물'은 사람에 대해 말해집니다. 따라서 `동물'은 개별 사람에 대해서도 말해집니다 --- 개별 사람은 사람이기도 하고 동물이기도 합니다. 이것은 류(\gr{γένος}) 술어의 전이성(transitivity)입니다. 현대 논리학의 용어로 말하면, said-of 관계는 전이적(transitive)이라는 것입니다. 이 원리는 나중에 5장에서 `다른 모든 것이 제1실체에 대해 말해지거나 제1실체 안에 있다'는 주장의 근거가 됩니다.

3장의 후반부(1b16--24)는 서로 종속되지 않는 류(\gr{γένος})들의 종차(\gr{διαφορά})가 서로 다르다는 원리를 제시합니다. 동물과 지식은 서로 종속되지 않는 류이므로, 동물의 종차(두 발 달린, 날개 달린, 수생의)는 지식의 종차가 아닙니다. 다만, 종속 관계에 있는 류들은 같은 종차를 가질 수 있습니다. 이 원리는 나중에 포르피리오스의 이사고게에서 `포르피리오스의 나무'(\gri{Arbor Porphyriana})로 시각화되는 류-종-종차의 위계 구조의 기초가 됩니다.

\subsection{열 가지 범주 (4장, 1b25--2a10)}
\vspace{-6pt}

\subsubsection{범주의 열거}

4장에서 아리스토텔레스는 결합 없이 말해지는 것들(2장에서 도입한 비결합 표현)을 열 가지 범주로 분류합니다. 이 열거는 서양 철학사에서 가장 유명한 분류 가운데 하나입니다.
\par\nobreak\vspace{-\parskip}

\srcquote{\gr{τῶν κατὰ μηδεμίαν συμπλοκὴν λεγομένων ἕκαστον ἤτοι οὐσίαν σημαίνει ἢ ποσὸν ἢ ποιὸν ἢ πρός τι ἢ ποῦ ἢ ποτὲ ἢ κεῖσθαι ἢ ἔχειν ἢ ποιεῖν ἢ πάσχειν.}}{`결합 없이 말해지는 것들은 각각 실체(\gr{οὐσία}), 양(\gr{ποσόν}), 질(\gr{ποιόν}), 관계(\gr{πρός τι}), 장소(\gr{ποῦ}), 시간(\gr{ποτέ}), 상태(\gr{κεῖσθαι}), 소유(\gr{ἔχειν}), 행위(\gr{ποιεῖν}), 겪음(\gr{πάσχειν})을 의미한다.'}{범주론 1b25--27}

아리스토텔레스는 이어서 각 범주의 예시를 제시합니다(1b27--2a4).

{%
\footnotesize
\setlength{\extrarowheight}{1.5pt}%
\renewcommand{\arraystretch}{1.1}%
\setlength{\tabcolsep}{3pt}%
\rowcolors{2}{altrowbg}{white}%
\noindent\centering
\begin{tabular}{|>{\centering\arraybackslash}m{5mm}|>{\centering\arraybackslash}m{11mm}|>{\centering\arraybackslash}m{17mm}|>{\centering\arraybackslash}m{48mm}|}
\hline
\rowcolor{tableheadblue}
{\color{h1navy}\bfseries \#} &
{\color{h1navy}\bfseries 범주} &
{\color{h1navy}\bfseries 그리스어} &
{\color{h1navy}\bfseries 아리스토텔레스의 예시} \\ \hline
1 & 실체 & \gr{οὐσία} & 사람, 말 \\ \hline
2 & 양 & \gr{ποσόν} & 두 자, 세 자 \\ \hline
3 & 질 & \gr{ποιόν} & 흰, 문법적인 \\ \hline
4 & 관계 & \gr{πρός τι} & 두 배, 절반, 더 큰 \\ \hline
5 & 장소 & \gr{ποῦ} & 뤼케이온에서, 광장에서 \\ \hline
6 & 시간 & \gr{ποτέ} & 어제, 작년에 \\ \hline
7 & 상태 & \gr{κεῖσθαι} & 누워 있다, 앉아 있다 \\ \hline
8 & 소유 & \gr{ἔχειν} & 신발을 신고 있다 \\ \hline
9 & 행위 & \gr{ποιεῖν} & 자르다, 태우다 \\ \hline
10 & 겪음 & \gr{πάσχειν} & 잘리다, 태워지다 \\ \hline
\end{tabular}
\par
}

\subsubsection{범주의 성격에 관한 물음: 왜 정확히 열 가지인가}

아리스토텔레스는 이 열 가지를 나열하지만, 왜 정확히 이 열 가지인지, 어떤 원리에서 이 목록이 도출되는지를 설명하지 않습니다. 이것은 고대부터 현대까지 범주론 해석의 가장 집요한 난제 가운데 하나입니다. 범주의 수가 다른 저작에서 달라지기도 합니다 --- 토피카에서는 일부가 빠지고, 분석론 후서에서는 다르게 열거됩니다.

고대 주석가들은 이 문제를 해결하려고 시도했습니다. 암모니오스(Ammonius)는 첫 네 범주(실체, 양, 질, 관계)를 결합하면 나머지 여섯이 도출된다고 주장했지만, 이 시도는 설득력이 약했습니다. 암모니오스의 제자 심플리키오스(Simplicius)는 더 정교한 도출을 시도했습니다.\footnote{Simplicius의 시도에 대해서는 Studtmann, ``Aristotle's Categories,'' §4 참조. 포르피리오스는 최소 분할이 4개이고 최대 분할이 10개라고만 말했을 뿐 도출을 시도하지 않았습니다(Porph.\ \gri{In Cat.}\ 71.20).} 근대에는 아돌프 트렌델렌부르크(Adolf Trendelenburg)가 1846년 대저 『범주론의 역사』(\gri{Geschichte der Kategorienlehre})에서 이 문제를 체계적으로 다루었고, 최근에는 스터트만이 아리스토텔레스의 질(quality) 범주에 대한 체계적 도출을 시도하면서, 이것이 범주 체계 전체의 체계적 도출 가능성에 대한 증거가 될 수 있다고 주장했습니다.\footnote{Paul Studtmann, ``Aristotle's Category of Quality: A Regimented Interpretation,'' \gri{Apeiron} 36, no.\ 3 (2003): 205--227; idem, ``Aristotle's Categorial Scheme,'' in \gri{The Oxford Handbook of Aristotle}, ed.\ Christopher Shields (Oxford: Oxford University Press, 2012), 63--80.}

한 가지 더 생각해 볼 점은, 이 열 가지 범주가 존재론적 분류(사물의 종류)인지 의미론적 분류(말해지는 것의 종류)인지 아니면 양자의 결합인지입니다. 1b25의 원문은 `의미한다'(\gr{σημαίνει})라는 동사를 사용하여 의미론적 색채를 띠지만, 5장 이하에서 아리스토텔레스는 범주를 존재론적 구조로 다룹니다. 이 긴장은 범주론 해석의 핵심 쟁점 가운데 하나이며, J.W.\ 소프(J.W.\ Thorp)는 이 문제를 정면으로 다루어 아리스토텔레스의 범주가 순수한 언어적 분류도, 순수한 존재론적 분류도 아닌 양자의 결합이라고 주장했습니다.\footnote{J.W.\ Thorp, ``Aristotle's Use of Categories,'' \gri{Phronesis} 19 (1974): 238--256.} 4장의 마지막(2a4--10)에서 아리스토텔레스는 비결합 표현 자체는 참도 거짓도 아니며, 결합에 의해서만 참·거짓이 성립한다고 재확인합니다.

\subsection{실체: 제1실체와 제2실체 (5장 전반, 2a11--4a21)}

\subsubsection{제1실체의 정의}

5장은 범주론에서 가장 길고 가장 중요한 장입니다. 아리스토텔레스는 실체(\gr{οὐσία}) 범주를 다른 아홉 범주와 구별되는 특별한 지위를 가진 것으로 논합니다. 장의 서두에서 제1실체의 정의가 제시됩니다.

\srcquote{\gr{Οὐσία δέ ἐστιν ἡ κυριώτατά τε καὶ πρώτως καὶ μάλιστα λεγομένη, ἣ μήτε καθ᾽ ὑποκειμένου τινὸς λέγεται μήτ᾽ ἐν ὑποκειμένῳ τινί ἐστιν, οἷον ὁ τὶς ἄνθρωπος ἢ ὁ τὶς ἵππος.}}{`가장 엄밀하게, 으뜸으로, 가장 탁월하게 실체라 불리는 것은, 어떤 주어에 대해서도 말해지지 않고 어떤 주어 안에도 있지 않은 것이다. 예컨대 개별 사람이나 개별 말이다.'}{범주론 2a11--14}

세 개의 최상급(\gr{κυριώτατά}, \gr{πρώτως}, \gr{μάλιστα})이 겹쳐 사용된 것에 주목해야 합니다. 아리스토텔레스는 `실체'라는 이름이 여러 수준에서 적용될 수 있음을 인정하면서, 그 가운데 가장 본래적인 의미의 실체가 바로 개별자라고 선언합니다. 이것은 사중 분류의 (d)칸을 존재론의 토대로 명시적으로 선언한 것입니다.

\subsubsection{제2실체}

제1실체가 속하는 종(\gr{εἶδος})과 그 종이 속하는 류(\gr{γένος})는 `제2실체'(\gr{δευτέρα οὐσία})라 불립니다(2a14--18). 예컨대 개별 사람은 `사람'이라는 종에 속하고, `사람'은 `동물'이라는 류에 속합니다. `사람'과 `동물' 모두 제2실체입니다. 제2실체 가운데서도 종이 류보다 더 실체적입니다 --- `이 사람은 무엇인가?'라고 물었을 때, `사람'이라고 답하는 것이 `동물'이라고 답하는 것보다 더 적합하기 때문입니다(2b7--14).

\subsubsection{제1실체의 존재론적 우선성}

5장에서 가장 핵심적인 논변은 2a34--2b6에서 전개됩니다. 아리스토텔레스는 다른 모든 것이 제1실체에 대해 말해지거나 제1실체 안에 있다고 주장하고, 여기서 결정적인 결론을 이끌어냅니다.

\srcquote{\gr{ὥστε τὰ ἄλλα πάντα ἤτοι καθ᾽ ὑποκειμένων τῶν πρώτων οὐσιῶν λέγεται ἢ ἐν ὑποκειμέναις αὐταῖς ἐστίν. εἰ οὖν αἱ πρῶται οὐσίαι μὴ ὑπῆρχον, ἀδύνατον τῶν ἄλλων τι εἶναι.}}{`따라서 다른 모든 것은 제1실체에 대해 말해지거나 제1실체 안에 있다. 그러므로 제1실체가 존재하지 않으면 다른 어떤 것도 존재할 수 없을 것이다.'}{범주론 2b4--6}

이것은 범주론의 존재론적 핵심입니다. 보편자(종과 류)도, 속성(질·양 등)도, 모두 궁극적으로 제1실체에 의존합니다.

% ── 유형 5: 제1실체 중심 의존 구조 방사 다이어그램 ──
\par\nobreak\vspace{6pt}
\begin{center}
\begin{tikzpicture}[
  scale=0.82, every node/.style={transform shape},
  box/.style={draw, rounded corners=2pt, minimum width=28mm, minimum height=8mm, align=center, font=\footnotesize},
  centerbox/.style={draw, rounded corners=2pt, minimum width=32mm, minimum height=10mm, align=center, fill=tableheadblue, font=\footnotesize\bfseries},
  arr/.style={-{Stealth[length=4pt]}, thick, color=midgray}
]
  % 중심: 제1실체
  \node[centerbox] (P) {{\color{h1navy}제1실체} \\ {\color{h1navy}\gr{πρώτη οὐσία}}};

  % 위: 제2실체
  \node[box, above=16mm of P] (S) {(a) 제2실체 \\ 예: 사람, 동물};

  % 왼쪽 아래: 개별적 속성
  \node[box, below left=14mm and 10mm of P] (IA) {(b) 개별적 속성 \\ 예: 이 흼, 이 문법 지식};

  % 오른쪽 아래: 보편적 속성
  \node[box, below right=14mm and 10mm of P] (UA) {(c) 보편적 속성 \\ 예: 지식, 색};

  % 화살표: 의존 방향 (→ 제1실체)
  \draw[arr] (S) -- node[right, font=\scriptsize, color=midgray] {말해짐} (P);
  \draw[arr] (IA) -- node[left, font=\scriptsize, color=midgray] {안에 있음} (P);
  \draw[arr] (UA) -- node[right, font=\scriptsize, color=midgray, align=left] {말해짐 \&\\ 안에 있음} (P);
\end{tikzpicture}

{\small\color{midgray} 제1실체의 존재론적 우선성: 다른 모든 존재자가 제1실체에 의존한다 (범주론 2b4--6).}
\end{center}
\vspace{-2pt}

이 존재론적 우선성을 시각적으로 표현하면, 제1실체가 중심에 놓이고 나머지 세 종류의 존재자가 모두 제1실체를 향해 의존 관계의 화살표를 보내는 방사 구조가 됩니다. 제2실체는 제1실체에 대해 말해지는 것(said-of)이고, 개별적 속성은 제1실체 안에 있는 것(present-in)이며, 보편적 속성은 양쪽 관계를 모두 맺습니다. 모든 화살표가 제1실체를 향한다는 것이 바로 2b4--6의 논변이 말하는 바입니다.

스탠퍼드 철학 백과사전의 「실체」 항목은, 범주론의 실체론이 중요한 논리적 구별들을 설정하지만 실체 자체의 형이상학적 분석에는 들어가지 않는다고 지적합니다. 그 분석은 주로 형이상학 Z권에서 이루어집니다.\footnote{SEP, ``Substance'' (Howard Robinson), §2.1.} 이 점은 범주론을 읽을 때 중요한 맥락입니다 --- 범주론은 존재론의 `지도'를 그리는 작업이지, 존재의 `본성'을 탐구하는 작업은 아닙니다.

\subsubsection{범주론과 형이상학의 긴장: `두 체계' 문제}

범주론과 형이상학의 관계는 아리스토텔레스 연구에서 가장 논쟁적인 주제 가운데 하나입니다. 범주론에서 `제1실체'는 개별자 --- 이 사람, 이 말 --- 이지만, 형이상학 Z권에서 아리스토텔레스는 실체의 으뜸가는 후보로 형상(\gr{μορφή}/\gr{εἶδος})을 선택하는 것으로 보입니다(형이상학 1028b35 이하). 이것은 데니얼 그레이엄(Daniel Graham)이 1987년 저서에서 `두 체계'(two systems)라고 부른 문제를 제기합니다.\footnote{Daniel Graham, \gri{Aristotle's Two Systems} (Oxford: Clarendon Press, 1987).}

학자들의 반응은 크게 세 갈래로 나뉩니다. 첫째, 그레이엄처럼 두 체계가 양립 불가능하며 아리스토텔레스의 사유가 발전적으로 변화했다고 보는 입장. 둘째, 범주론의 목적이 단순히 형이상학적 분석을 요구하지 않았을 뿐이며, 두 저작이 같은 존재론의 다른 측면을 다룬다고 보는 입장. 셋째, 데이비드 위긴스(David Wiggins)처럼 범주론에서 형이상학으로의 변화가 오히려 해롭다고 보는 입장입니다. 위긴스는 형이상학에서 형상을 실체로 선택한 것이 당혹스럽다고 평하는데, 형상이 보편자라면 제2실체에 해당하므로 범주론의 기준에서는 제1실체가 될 수 없기 때문입니다.\footnote{David Wiggins, \gri{Sameness and Substance Renewed} (Cambridge: Cambridge University Press, 1998), 232 이하. 다만, 아리스토텔레스의 실체적 형상(substantial form)이 정말로 보편자인지는 그 자체가 논쟁적인 문제입니다.}

이 논쟁은 42주 강독의 범위를 넘지만, 범주론을 읽을 때 이 긴장을 의식하는 것은 중요합니다.

\subsubsection{실체의 특성들}

5장의 후반부(3a7--4a21)는 실체에 공통되는 여러 특성을 열거합니다.

첫째, 실체는 어떤 주어 안에도 있지 않습니다(3a7--21). 이것은 제1실체에는 자명하지만, 제2실체에도 해당됩니다 --- `사람'은 개별 사람에 대해 말해지지만 개별 사람 `안에' 있는 것은 아닙니다. 아리스토텔레스는 여기서 한 가지 예외를 인정합니다: 종차(\gr{διαφορά})도 주어 안에 있지 않습니다(3a21--28).

둘째, 실체와 종차에서 파생되는 술어는 모두 동의적(\gr{συνωνύμως})으로 술어됩니다(3a33--b9). 여기서 1장의 동음이의·동의 구별이 다시 작동합니다. `사람'이 소크라테스에 대해 술어될 때, 이름뿐 아니라 정의도 함께 전이됩니다. 반면 `흰'이 소크라테스에 대해 술어될 때는 이름은 전이되지만 정의는 전이되지 않을 수 있습니다. 이것이 바로 said-of 관계와 present-in 관계의 차이가 술어화의 차원에서 나타나는 방식입니다.

셋째, 모든 실체는 `이것'(\gr{τόδε τι})을 의미하는 것으로 보입니다(3b10--23). 이 대목은 아리스토텔레스 존재론에서 핵심적인 개념을 도입합니다.

\srcquote{\gr{Πᾶσα δὲ οὐσία δοκεῖ τόδε τι σημαίνειν. ἐπὶ μὲν οὖν τῶν πρώτων οὐσιῶν ἀναμφισβήτητον καὶ ἀληθές ἐστιν ὅτι τόδε τι σημαίνει· ἄτομον γὰρ καὶ ἓν ἀριθμῷ τὸ δηλούμενόν ἐστιν.}}{`모든 실체는 «이것»(\gr{τόδε τι})을 의미하는 것으로 보인다. 제1실체에 대해서는 이것이 논쟁의 여지 없이 참이다 --- 드러나는 것이 분할 불가능하고(\gr{ἄτομον}) 수적으로 하나(\gr{ἓν ἀριθμῷ})이기 때문이다.'}{범주론 3b10--13}

그러나 아리스토텔레스는 곧바로 이것이 제2실체에게는 엄밀히 참이 아니라고 지적합니다. 제2실체는 오히려 `어떤 종류의 것'(\gr{ποιόν τι})을 의미합니다 --- 다만 질(\gr{ποιόν})처럼 단순한 성질을 의미하는 것이 아니라, `실체의 한정'(\gr{περὶ οὐσίαν τὸ ποιὸν ἀφορίζει})을 의미합니다.

넷째, 실체에는 반대(\gr{ἐναντίον})가 없습니다(3b24--32).

다섯째, 실체는 정도의 차이(\gr{μᾶλλον καὶ ἧττον})를 허용하지 않습니다(3b33--4a9).

\subsubsection{실체의 가장 고유한 특성: 반대자를 수용하는 능력}

5장의 마지막 논변(4a10--4a21, 그리고 이후 4b19까지 이어지나 이번 주 범위는 4a21까지)은 실체의 가장 고유한(\gr{μάλιστα ἴδιον}) 특성을 제시합니다. 수적으로 하나이고 동일한 것이면서 반대자들(\gr{τὰ ἐναντία})을 수용할 수 있다는 것입니다. 예컨대 수적으로 하나인 색은 동시에 검고 희지 않으며, 수적으로 하나인 행위가 동시에 나쁘고 좋지도 않습니다. 그러나 수적으로 하나인 실체 --- 예컨대 개별 사람 --- 는 어떤 때에는 희고 어떤 때에는 검으며, 어떤 때에는 뜨겁고 어떤 때에는 차갑습니다.

이 특성이 왜 `가장 고유한' 것인지를 생각해 보면, 앞의 네 특성은 모두 실체만의 것이 아니었습니다. 반대자를 수용하는 능력만이 실체에 고유합니다. 그리고 이 능력은 제1실체의 존재론적 우선성과 직결됩니다 --- 변화의 주어(\gr{ὑποκείμενον})로서 동일성을 유지하면서 반대되는 속성들을 번갈아 받아들이는 것은 바로 개별적 실체만이 할 수 있는 일이기 때문입니다. 스탠퍼드 철학 백과사전의 「실체」 항목이 지적하듯이, \citework{Phys}{}에서도 아리스토텔레스는 모든 변화가 기저하는 주어(\gr{ὑποκείμενον})를 전제한다고 논변하며, 이 주어가 바로 실체입니다.\footnote{SEP, ``Substance,'' §2.1. `실체'라는 이름 자체가 `아래에 놓여 있는 것'(\gr{ὑποκείμενον}, \gri{sub-stantia})이라는 뜻을 가진다는 점에 주목할 가치가 있습니다 --- 라틴어 \gri{substantia}는 \gri{sub}(아래에) + \gri{stare}(서 있다)의 합성어입니다.}

다음 주(3주차)에 이 논변의 나머지와, `진술과 믿음도 반대자를 수용하지 않는가?'라는 반론(4a22--4b19)을 검토하겠습니다.

\begin{references}
\subsubsection*{참고문헌}

\refitem Ackrill, J.\ L. \gri{Aristotle's Categories and De Interpretatione}. Clarendon Aristotle Series. Oxford: Clarendon Press, 1963.

\refitem Cohen, Marc. ``Predication and Ontology: Categories.'' Lecture notes, University of Washington.

\refitem Frede, Michael. ``Categories in Aristotle.'' In \gri{Essays in Ancient Philosophy}, 29--48. Minneapolis: University of Minnesota Press, 1987.

\refitem Graham, Daniel. \gri{Aristotle's Two Systems}. Oxford: Clarendon Press, 1987.

\refitem Hintikka, Jaakko. ``Semantical Games, the Alleged Ambiguity of `Is' and Aristotelian Categories.'' \gri{Synthese} 54 (1983): 443--468.

\refitem Minio-Paluello, L., ed. \gri{Aristotelis Categoriae et Liber de Interpretatione}. OCT. Oxford: Clarendon Press, 1949; repr.\ with corrections, 1956.

\refitem Owen, G.E.L. ``Logic and Metaphysics in Some Earlier Works of Aristotle.'' In \gri{Aristotle and Plato in the Mid-Fourth Century}, edited by I.\ Düring and G.E.L.\ Owen, 163--190. Göteborg, 1960.

\refitem Robinson, Howard. ``Substance.'' In \gri{Stanford Encyclopedia of Philosophy}.\newline https://plato.stanford.edu/entries/substance/.

\refitem Studtmann, Paul. ``Aristotle's Categories.'' In \gri{Stanford Encyclopedia of Philosophy}. https://plato.stanford.edu/entries/aristotle-categories/.

\refitem ---. ``Aristotle's Category of Quality: A Regimented Interpretation.'' \gri{Apeiron} 36, no.\ 3 (2003): 205--227.

\refitem ---. ``Aristotle's Categorial Scheme.'' In \gri{The Oxford Handbook of Aristotle}, edited by Christopher Shields, 63--80. Oxford: Oxford University Press, 2012.

\refitem Thorp, J.\ W. ``Aristotle's Use of Categories.'' \gri{Phronesis} 19 (1974): 238--256.

\refitem Trendelenburg, Adolf. \gri{Geschichte der Kategorienlehre}. Berlin: Verlag von G.\ Bethge, 1846.

\refitem Wedin, Michael V. ``Nonsubstantial Individuals.'' \gri{Phronesis} 38 (1993): 137--165.

\refitem Wiggins, David. \gri{Sameness and Substance Renewed}. Cambridge: Cambridge University Press, 1998.

\end{references}

\end{document}
